\section{The variance inequality}
\label{sec:variance}

We now apply the moment-coordinate estimates to the potential
\eqref{eq:potential}. In this section all expectations are with
respect to its probability $\mu$.

\begin{proof}[Proof of Proposition~\ref{prop:variance}]
The potential is smooth and even, and $D^2V\ge\eps I$.
Lemma~\ref{lem:moment} therefore supplies the moment coordinates,
with $0<\tau\le (4/\eps)I$. Put
\[
 H_i=D_x^2\psi_i,\qquad a_i=\E\psi_i.
\]
Each $H_i$ is positive semidefinite, and $a_i>0$ since $p_i\ne0$
and $\mu$ has positive density everywhere. Homogeneity gives
\begin{equation}\label{eq:homogeneity}
 x\cdot\nabla\psi_i=q\psi_i,
 \qquad H_ix=(q-1)\nabla\psi_i.
\end{equation}
The adjunction identity of Section~\ref{sec:identities} also gives
\begin{equation}\label{eq:mi}
 m_i:=\E\tr(\tau H_i)=\E x\cdot\nabla\psi_i=q a_i.
\end{equation}
To justify it for these noncompact functions, apply
$\E x_j k=\E\tau_{jk}\partial_{x_k}k$ to
$k=\eta_R\partial_{x_j}\psi_i$ and sum in $j$, where $\eta_R$
is a smooth cutoff equal to one on $B_R$ and zero outside $B_{2R}$.
The extra term is bounded in absolute value by
$C R^{-1}\E[|\nabla\psi_i|\mathbf1_{\{|x|\ge R\}}]$, which tends
to zero. The other terms converge because $\tau$ is bounded and
all polynomial moments are finite.

Take an even $f\in C_c^\infty(\R^n)$ and use
\[
 B=\tau D_x^2f\,\tau,\qquad W=\delta_\mu B,
 \qquad u=\delta_\mu W=A(A-1)f,
 \qquad M=\tau^{-1}.
\]
Applying \eqref{eq:test-pairing} to $F=\psi_i$ gives, for each $i$,
\begin{equation}\label{eq:di}
 d_i:=\E u\psi_i=\E W\cdot\nabla\psi_i=\E\tr(BH_i).
\end{equation}
Two Cauchy--Schwarz estimates will give complementary bounds for
these same numbers $d_i$.

First apply the Hilbert--Schmidt Cauchy--Schwarz inequality, including
integration, to the matrix fields
$H_i^{1/2}\tau^{1/2}$ and $H_i^{1/2}B\tau^{-1/2}$. Their pairing
is $\tr(H_iB)$, so
\begin{equation}\label{eq:first-cs}
 d_i^2\le m_i\,\E\tr(MBH_iB).
\end{equation}
This does not require $B$ to be positive semidefinite.
Second, \eqref{eq:homogeneity}--\eqref{eq:di} give
\[
 \E x^TH_i x=(q-1)m_i,
 \qquad \E x^TH_iW=(q-1)d_i.
\]
Cauchy--Schwarz for the positive semidefinite form
$(X,Y)\mapsto\E X^TH_iY$ therefore yields
\begin{equation}\label{eq:second-cs}
 \E W^TH_iW\ge(q-1)\frac{d_i^2}{m_i}.
\end{equation}

Since $H=D^2V=\eps I+\sum_iH_i$, Lemma~\ref{lem:tensor} and
the weighted Bochner identity \eqref{eq:bochner} imply
\begin{align}
 \E u^2
 &\ge\E\tr(MBHB)+\E W^THW\notag\\
 &\ge\sum_i\left(\E\tr(MBH_iB)+\E W^TH_iW\right)
 \ge\sum_i\frac{q d_i^2}{m_i}
 =\sum_i\frac{d_i^2}{a_i}.
 \label{eq:dual-bound}
\end{align}
The discarded terms are nonnegative, since $M$ is positive definite
and $B$ is symmetric. This is the required estimate on the dense
class of double divergences.

By Lemma~\ref{lem:density}, such $u$ are dense in the centered even
subspace of $L^2(\mu)$. Each $\psi_i\in L^2(\mu)$, so all pairings
$u\mapsto\E u\psi_i$ are continuous. Thus \eqref{eq:dual-bound}
holds for every centered even $u\in L^2(\mu)$.
For $F=\sum_i b_i\psi_i$, choose $u=F-\E F$. Then
\[
 \|u\|_2^2=\sum_i b_i d_i
 \le\left(\sum_i b_i^2a_i\right)^{1/2}
      \left(\sum_i d_i^2/a_i\right)^{1/2}
 \le\left(\sum_i b_i^2a_i\right)^{1/2}\|u\|_2.
\]
If $\|u\|_2=0$, the conclusion is immediate; otherwise division and
squaring prove \eqref{eq:variance}. This completes the analytic
assertion and, by Section~\ref{sec:reduction}, Theorem~\ref{thm:main}.
\end{proof}
